\clearpage
\fancyhead[L]{\small HARDING. HIMMELSKARTEN.}
\hypertarget{native-598-lower}{}
\begin{center}\small Göttingische gelehrte Anzeigen. Stück 5. Seiten 41..45. 9. Januar 1817.\\[3pt]\rule{.72\linewidth}{.3pt}\end{center}

Von den Himmelskarten unseres Herrn Prof.\ \textsc{Harding}, deren zuletzt erschienene Lieferungen wir in Stück 158 dieser Zeitschrift von 1812 [S. 541 d. B.] angezeigt haben, ist jetzt wieder im Kommissionsvertrieb der Vandenhoeck-Ruprechtschen Buchhandlung eine neue, die \emph{fünfte}, Lieferung mit vier Blättern erschienen, nämlich die Nummern VI, X, XIX und XXV. Damit man leichter überblicken kann, wie viel von dieser nützlichen Arbeit jetzt vollendet ist und wie viel noch fehlt, erinnern wir hier nur kurz an ihren Anlass und ihren Plan. Die Entdeckung der Planeten Ceres und Pallas, die nie größer als Fixsterne achter Größenklasse und oft noch viel kleiner erscheinen, machte erstmals für die Teile des Himmels, in denen jene interessanten Himmelskörper sich zeigen können, weit detailliertere Sternkarten unentbehrlich, als wir bis dahin besaßen. Um nicht jedes Jahr besondere Karten für sie zu benötigen, war ein eigener Atlas für den Tierkreis dieser Planeten wünschenswert, den der Verfasser zu erarbeiten beschloss. Eine mäßige Anzahl von Blättern (10 oder 11) hätte dafür ausgereicht. Doch während der Verfasser diese Arbeit bereits verfolgte, kam, zum Teil gerade durch diese Arbeit veranlasst, die Entdeckung zweier weiterer neuer Planeten hinzu, die eine Erweiterung des Plans notwendig machte. Der vielfältige Nutzen, den detaillierte Stern\SourceLeaf{599}karten auch in anderer Hinsicht bieten, bewog den Verfasser, seinen Atlas, von dem inzwischen schon einige Blätter fertiggestellt waren, auf den gesamten nördlichen Himmel und auf den Teil des südlichen Himmels auszudehnen, der in unseren Gegenden Europas noch gut beobachtet werden kann und für den noch die Materialien vorhanden sind, also bis etwas mehr als $30^\circ$ südlicher Deklination. Neun Blätter mit einer Ausdehnung von jeweils $40^\circ$ in Rektaszension sollten also diesen südlichen Teil des Himmels abdecken. Ebenso sollten neun Blätter von jeweils ähnlicher Ausdehnung den Teil des nördlichen Himmels vom Äquator bis zu $32^\circ$ nördlicher Deklination abdecken. Diese 18 Blätter sind nun alle fertiggestellt, nachdem die bisher noch fehlenden Blätter VI und X in der vorliegenden Lieferung hinzugekommen sind. Allerdings wird künftig Blatt Nr.\ I, das zu Beginn noch nach dem begrenzten Plan ausgearbeitet wurde, sich von beinahe $22^\circ$ südlicher bis fast $13^\circ$ nördlicher Deklination erstreckt und dessen nördlicher Teil auch auf dem neuen Blatt Nr.\ X enthalten ist, aus dem Atlas genommen und durch ein anderes ersetzt, das bis $32^\circ$ südlicher Deklination reicht und den Besitzern des Atlas demnächst vom Verfasser unentgeltlich nachgeliefert werden wird.

Der Himmelsbereich, den diese 18 Blätter umfassen, ist für die Astronomen insofern am interessantesten, als sämtliche Planeten und der Mond nie über ihn hinausgehen. Deshalb hielt der Verfasser es mit Recht für geboten, zuerst diesen Bereich fertigzustellen. Eine ebenso detaillierte Bearbeitung des nördlichen Himmels von $32^\circ$ Deklination bis zum Pol wäre allerdings schon wegen der Kometen, die ihn so häufig durchstreifen, als wesentliches Bedürfnis anzusehen. Denn viele Kometen sind \emph{nur} im Fernrohr sichtbar, und selbst die größeren lassen sich noch eine Zeit lang mit Fernrohren verfolgen, wenn ihre zunehmende Entfernung sie dem bloßen Auge bereits entzogen hat. Bei solchen Gelegenheiten muss sich der Beobachter in Ermangelung gestochener Sternkarten stets der zeitraubenden Arbeit unterziehen, selbst eine Karte von der Gegend zu entwerfen, die der Komet durchläuft. Die Blätter XIX und XXV des vorliegenden Atlas machen den Anfang, diesem Bedürfnis abzuhelfen. Der dabei befolgte Plan zur Aufteilung des nördlichen Himmels ist der folgende. Da die Grade auf den Parallelkreisen zum Pol hin immer kleiner werden, konnte ein Blatt von derselben Größe wie die ersten 18 einen größeren Bereich in Rektaszension umfassen. Zuerst werden acht Blätter folgen, die sich jeweils am unteren Rand über etwas mehr als $45^\circ$ in Rektaszension und in der Mitte von $30^\circ$ bis fast $65^\circ$ Deklination erstrecken. Der übrige Teil des Himmels von $64^\circ$ Deklination bis zum Pol lässt sich dann bei nur wenig verkleinertem Maßstab auf zwei Blättern unterbringen. Die Ecken, in denen sonst bereits auf den Blättern gezeigte Himmelsbereiche wiederholt werden müssten, sollen zur Darstellung interessanter Sterngruppen wie der Plejaden, der Krippe und anderer dienen. Mit 28 Blättern wird also der ganze Athas vollendet sein, und die jetzt noch fehlenden acht Blätter, die größtenteils schon gezeichnet sind, werden sich hoffentlich bald liefern lassen.

Bei den ersten 18 Blättern war eine aufwendigere Projektionsart überflüssig. Das Netz aus einfachen Quadraten erleichterte die Zeichnung und bietet auch beim Gebrauch einige Annehmlichkeiten. Die geringe, ohnehin kaum merkliche Abweichung vom wahren Verhältnis der Rektaszensions- und Deklinationsgrade zum Rand hin fiel dabei überhaupt nicht ins Gewicht. Bei größeren Deklinationen war dies allerdings nicht mehr anwendbar. Für die acht Blätter XIX bis XXVI wählte der Verfasser daher eine ihm von Herrn Hofrat \textsc{Gauß} dafür vorgeschlagene Projektion, bei der die Deklinationskreise gerade Linien und die Parallelkreise konzentrische Kreise sind. Deren abnehmende Radien sind so genau bestimmt, dass die Rektaszensionsgrade überall im richtigen Verhältnis zu den Deklinationsgraden stehen. Bei dieser Darstellung ist das Bild dem Himmel also in den kleinsten Teilen vollkommen ähnlich, was die wesentlichste Bedingung ist, \SourceLeaf{600}Die Deklinationsgrade sind freilich, genau genommen, nicht alle gleich groß. Dieser in der Natur der Sache liegende unvermeidliche Umstand spielt hier jedoch keine Rolle, zumal der Unterschied so gering ist, dass er kaum bemerkt wird. Auf dem 64sten und 3osten Grad der Deklination misst ein Deklinationsgrad nämlich 13,07 Millimeter, in der Mitte auf dem 47sten Grad dagegen 12,38 Millimeter.

Da die Verteilung des Himmels auf die Karten nur durch die Rektaszensionen und Deklinationen bestimmt wird, erscheinen manche Sternbilder \mbox{naturgemäß} aufgeteilt. Blatt Nr.\ VI enthält Teile der Jungfrau, der Waage, des Skorpions, der Schlange, des Wolfs, des Zentauren und des Schwanzes der Hydra sowie den ganzen Vogel Einsiedler. Blatt Nr.\ X enthält Teile der Fische, der Andromeda, der Dreiecke, des Widders, des Walfischs und des Pegasus. Blatt Nr.\ XIX enthält Teile des Kepheus, der Kassiopeia, des Erntehüters, des Camelopard, des Perseus, der Andromeda, der Fische, der Dreiecke und des Widders. Blatt Nr.\ XXV enthält Teile des Drachen, des Herkules, der Leier, des Schwans und des Kepheus. Papier und Stich stehen an Schönheit den früheren Lieferungen nicht nach.
